% ============================================================
\chapter{Mod\`eles Stochastiques}
\label{ch:stochastiques}
% ============================================================

\begin{center}
\textit{<<\,Le hasard n'est que la mesure de notre ignorance.\,>>}\\[4pt]
--- Henri Poincar\'e
\end{center}

\bigskip

Jusqu'ici, nos mod\`eles \'etaient d\'eterministes~: connaissant l'\'etat initial et les lois d'\'evolution, on pouvait pr\'edire le futur avec certitude. Mais le monde r\'eel est bruit\'e. Les march\'es financiers fluctuent, les g\`enes mutent al\'eatoirement, les mol\'ecules s'agitent sous l'effet thermique, et m\^eme les \'epid\'emies ne se propagent pas de mani\`ere d\'eterministe. Mod\'eliser ces ph\'enom\`enes exige d'int\'egrer l'al\'ea dans les \'equations. Ce chapitre introduit les mod\`eles stochastiques --- des marches al\'eatoires aux \'equations diff\'erentielles stochastiques d'It\^o --- et montre comment le hasard, loin d'\^etre un obstacle \`a la compr\'ehension, devient un ingr\'edient essentiel de la mod\'elisation. Les mod\`eles stochastiques d\'ecrivent l'\'evolution de syst\`emes soumis
\`a des fluctuations al\'eatoires, ce qui est essentiel en biologie,
en th\'eorie des files d'attente et en physique.

% ------------------------------------------------------------
\section{Processus de naissance et de mort}
\label{sec:naissance_mort}
% ------------------------------------------------------------

\begin{definition}[Processus de naissance et de mort]
\index{processus de naissance et de mort}
Un \textbf{processus de naissance et de mort}\index{naissance et mort} est un processus
de Markov \`a temps continu $\{X(t)\}_{t \geq 0}$ \`a valeurs dans $\N$
dont les seules transitions possibles depuis l'\'etat $n$ sont~:
\begin{itemize}
  \item $n \to n+1$ (naissance) avec taux $\lambda_n \geq 0$,
  \item $n \to n-1$ (mort) avec taux $\mu_n \geq 0$, o\`u $\mu_0 = 0$.
\end{itemize}
\end{definition}

\begin{intuition}
Imaginez une population de bact\'eries dans une bo\^ite de P\'etri.
Chaque bact\'erie se divise (naissance) ou meurt ind\'ependamment des
autres. Le nombre total de bact\'eries fluctue al\'eatoirement autour
d'une tendance g\'en\'erale. Le processus de naissance et de mort
capture exactement cette dynamique.
\end{intuition}

\begin{proposition}[\'Equations de Kolmogorov]
\label{prop:kolmogorov_nm}
Soit $p_n(t) = \PP(X(t) = n)$. Les probabilit\'es d'\'etat satisfont le
syst\`eme d'\'equations diff\'erentielles~:
\[
  \frac{dp_n}{dt} = \lambda_{n-1}\, p_{n-1}(t) + \mu_{n+1}\, p_{n+1}(t)
    - (\lambda_n + \mu_n)\, p_n(t), \quad n \geq 1,
\]
avec $\dfrac{dp_0}{dt} = \mu_1\, p_1(t) - \lambda_0\, p_0(t)$.
\end{proposition}

\begin{proof}
On conditionne sur l'\'etat \`a l'instant $t$ et on utilise les propri\'et\'es
markoviennes. En un intervalle $[t, t+h]$, avec $h$ petit~:
\begin{align*}
  p_n(t+h) &= p_n(t)\bigl(1 - (\lambda_n + \mu_n)h\bigr)
    + p_{n-1}(t)\,\lambda_{n-1}\,h + p_{n+1}(t)\,\mu_{n+1}\,h + o(h).
\end{align*}
En soustrayant $p_n(t)$, divisant par $h$ et passant \`a la limite $h \to 0$,
on obtient l'\'equation annonc\'ee.
\end{proof}

\begin{theoreme}[Distribution stationnaire]
\label{thm:stat_nm}
Si les taux v\'erifient $\sum_{n=1}^{\infty} \prod_{k=0}^{n-1}
\frac{\lambda_k}{\mu_{k+1}} < \infty$, alors il existe une distribution
stationnaire $\pi = (\pi_n)_{n \geq 0}$ donn\'ee par~:
\[
  \pi_n = \pi_0 \prod_{k=0}^{n-1} \frac{\lambda_k}{\mu_{k+1}},
  \qquad \pi_0 = \left(1 + \sum_{n=1}^{\infty} \prod_{k=0}^{n-1}
  \frac{\lambda_k}{\mu_{k+1}}\right)^{-1}.
\]
\end{theoreme}

\begin{exemple}
Consid\'erons un processus de naissance pure avec $\lambda_n = \lambda\, n$
et $\mu_n = 0$ (mod\`ele de Yule\index{mod\`ele de Yule}).
Partant de $X(0)=1$, la population cro\^it exponentiellement~:
$\E[X(t)] = e^{\lambda t}$.
\end{exemple}

% ------------------------------------------------------------
\section{Th\'eorie des files d'attente}
\label{sec:files_attente}
% ------------------------------------------------------------

\begin{definition}[File M/M/1]
\index{file M/M/1}
Une file \textbf{M/M/1} est un processus de naissance et de mort avec
taux d'arriv\'ee constant $\lambda_n = \lambda$ et taux de service
constant $\mu_n = \mu$ pour $n \geq 1$. <<\,M\,>> d\'esigne
<<\,markovien\,>> (arriv\'ees et services exponentiels).
\end{definition}

\begin{theoreme}[\'Equilibre de la file M/M/1]
\label{thm:mm1}
Soit $\rho = \lambda / \mu$ l'intensit\'e du trafic\index{intensit\'e du trafic}.
Si $\rho < 1$, la file M/M/1 admet une distribution stationnaire
g\'eom\'etrique~:
\[
  \pi_n = (1-\rho)\,\rho^n, \quad n \geq 0.
\]
Le nombre moyen de clients dans le syst\`eme est
$L = \dfrac{\rho}{1-\rho}$
et le temps moyen de s\'ejour est
$W = \dfrac{1}{\mu - \lambda}$.
\end{theoreme}

\begin{proof}
Par le th\'eor\`eme~\ref{thm:stat_nm}, $\pi_n = \pi_0 \rho^n$.
La condition de normalisation $\sum_{n=0}^{\infty} \pi_n = 1$ donne
$\pi_0 = 1 - \rho$ lorsque $\rho < 1$. Le nombre moyen suit de
$L = \sum_{n=0}^{\infty} n\,\pi_n = (1-\rho) \sum_{n=1}^{\infty} n\,\rho^n
= \rho/(1-\rho)$. La formule de Little\index{formule de Little}
$L = \lambda W$ donne $W$.
\end{proof}

\begin{formules}
\textbf{R\'esum\'e des files d'attente}
\begin{align*}
  \text{M/M/1 :}\quad & L = \frac{\rho}{1-\rho}, \quad
    W = \frac{1}{\mu - \lambda}, \quad
    L_q = \frac{\rho^2}{1-\rho} \\[6pt]
  \text{M/M/c :}\quad & \pi_0 = \left[\sum_{k=0}^{c-1}
    \frac{(c\rho)^k}{k!} + \frac{(c\rho)^c}{c!\,(1-\rho)}\right]^{-1},
    \quad \rho = \frac{\lambda}{c\mu} < 1
\end{align*}
\end{formules}

\begin{definition}[File M/M/c]
\index{file M/M/c}
Une file \textbf{M/M/c} poss\`ede $c$ serveurs identiques en parall\`ele.
Les taux de service effectifs sont $\mu_n = \min(n, c)\,\mu$.
\end{definition}

\begin{remarque}
La formule d'Erlang~C\index{formule d'Erlang~C} donne la probabilit\'e
d'attente dans une file M/M/c~:
\[
  C(c, \lambda/\mu) = \frac{(c\rho)^c}{c!\,(1-\rho)} \cdot \pi_0.
\]
Elle est largement utilis\'ee dans le dimensionnement des centres d'appels.
\end{remarque}

% ------------------------------------------------------------
\section{Marches al\'eatoires}
\label{sec:marches}
% ------------------------------------------------------------

\begin{definition}[Marche al\'eatoire simple]
\index{marche al\'eatoire}
Une \textbf{marche al\'eatoire simple} sur $\Z$ est une suite
$(S_n)_{n \geq 0}$ d\'efinie par $S_0 = 0$ et $S_n = \sum_{k=1}^n X_k$,
o\`u les $(X_k)$ sont i.i.d. avec $\PP(X_k = +1) = p$ et $\PP(X_k = -1) = 1-p$.
\end{definition}

\begin{theoreme}[R\'ecurrence de la marche al\'eatoire]
\label{thm:recurrence_marche}
La marche al\'eatoire simple sur $\Z$ est r\'ecurrente si et seulement si $p = 1/2$.
Plus g\'en\'eralement, la marche al\'eatoire simple sur $\Z^d$ est r\'ecurrente
si et seulement si $d \leq 2$ (th\'eor\`eme de P\'olya\index{th\'eor\`eme de P\'olya}).
\end{theoreme}

\begin{proof}[Id\'ee de la preuve pour $d=1$]
Si $p = 1/2$, la probabilit\'e de retour \`a l'origine est
$\sum_{n=1}^{\infty} \PP(S_{2n} = 0) = \sum_{n=1}^{\infty}
\binom{2n}{n} 2^{-2n}$. Par la formule de Stirling,
$\binom{2n}{n} \sim \frac{4^n}{\sqrt{\pi n}}$, donc la s\'erie
diverge, ce qui assure la r\'ecurrence. Si $p \neq 1/2$,
$\E[X_k] = 2p-1 \neq 0$ et la loi des grands nombres montre
que $S_n \to \pm\infty$.
\end{proof}

\begin{exemple}
La marche al\'eatoire mod\'elise la position d'une particule
en suspension dans un fluide (mouvement brownien discret\index{mouvement brownien}).
Elle mod\'elise \'egalement la fortune d'un joueur dans un jeu \'equitable.
\end{exemple}

% ------------------------------------------------------------
\section{Algorithme de Gillespie}
\label{sec:gillespie}
% ------------------------------------------------------------

\begin{definition}[Algorithme de Gillespie (SSA)]
\index{algorithme de Gillespie}
L'\textbf{algorithme de simulation stochastique} (SSA) de Gillespie
est une m\'ethode exacte pour simuler la trajectoire d'un syst\`eme
chimique stochastique. \`A chaque \'etape, on d\'etermine~:
\begin{enumerate}
  \item le temps $\tau$ avant la prochaine r\'eaction (exponentiel
    de param\`etre $a_0 = \sum_j a_j$),
  \item la r\'eaction $j$ qui se produit (avec probabilit\'e $a_j / a_0$),
\end{enumerate}
o\`u $a_j$ est la \emph{propensit\'e} de la r\'eaction $j$.
\end{definition}

\begin{attention}
L'algorithme de Gillespie simule \emph{chaque} r\'eaction individuellement.
Pour des syst\`emes avec un grand nombre de mol\'ecules, il devient
tr\`es co\^uteux. On pr\'ef\`ere alors des m\'ethodes approch\'ees comme
le $\tau$-leaping\index{tau-leaping@$\tau$-leaping}.
\end{attention}

\begin{exemple}
Consid\'erons la r\'eaction de production-d\'egradation d'un ARN messager~:
$\emptyset \xrightarrow{k_1} \text{ARN} \xrightarrow{k_2} \emptyset$,
avec propensit\'es $a_1 = k_1$ et $a_2 = k_2\, n$ o\`u $n$ est le
nombre de mol\'ecules d'ARN. \`A l'\'equilibre stochastique, $n$ suit
une loi de Poisson de param\`etre $k_1/k_2$.
\end{exemple}

\begin{minted}{python}
import numpy as np

def gillespie_production_degradation(k1, k2, n0, t_max):
    """Simulation de Gillespie pour production-degradation."""
    t, n = 0.0, n0
    times, states = [t], [n]
    while t < t_max:
        a1 = k1           # propension production
        a2 = k2 * n       # propension degradation
        a0 = a1 + a2
        if a0 == 0:
            break
        tau = np.random.exponential(1.0 / a0)
        t += tau
        if np.random.uniform() < a1 / a0:
            n += 1         # production
        else:
            n -= 1         # degradation
        times.append(t)
        states.append(n)
    return np.array(times), np.array(states)
\end{minted}

% ------------------------------------------------------------
\section{Simulation de Monte Carlo}
\label{sec:monte_carlo}
% ------------------------------------------------------------

\begin{definition}[M\'ethode de Monte Carlo]
\index{Monte Carlo}
Une \textbf{m\'ethode de Monte Carlo} est une technique de calcul qui
utilise des \'echantillons al\'eatoires pour estimer une quantit\'e
d\'eterministe. Pour estimer $\theta = \E[g(X)]$, on g\'en\`ere
$X_1, \ldots, X_N$ i.i.d. selon la loi de $X$ et on calcule~:
\[
  \hat{\theta}_N = \frac{1}{N} \sum_{i=1}^N g(X_i).
\]
\end{definition}

\begin{theoreme}[Convergence de Monte Carlo]
Par la loi forte des grands nombres, $\hat{\theta}_N \to \theta$ p.s.
De plus, par le TCL, l'erreur d'estimation est d'ordre $O(1/\sqrt{N})$~:
\[
  \sqrt{N}\,(\hat{\theta}_N - \theta) \xrightarrow{\mathcal{L}}
  \mathcal{N}(0, \Var(g(X))).
\]
\end{theoreme}

\begin{intuition}
La convergence en $O(1/\sqrt{N})$ est \emph{ind\'ependante de la dimension}
du probl\`eme. C'est l'avantage majeur de Monte Carlo par rapport
aux m\'ethodes d\'eterministes (quadrature), dont la convergence se
d\'egrade avec la dimension (<<\,fl\'eau de la dimension\,>>).
\end{intuition}

\begin{proposition}[R\'eduction de variance par variables antith\'etiques]
\index{variables antith\'etiques}
Si $g$ est monotone et $X \sim \mathcal{U}(0,1)$, l'estimateur
\[
  \hat{\theta}_N^{\text{anti}} = \frac{1}{N}\sum_{i=1}^{N/2}
  \frac{g(U_i) + g(1-U_i)}{2}
\]
a une variance inf\'erieure ou \'egale \`a celle de l'estimateur brut.
\end{proposition}

\begin{exercice}
Estimer $\int_0^1 e^x\,dx$ par Monte Carlo avec $N = 10\,000$ \'echantillons.
Comparer la pr\'ecision obtenue avec et sans variables antith\'etiques.
\end{exercice}

\begin{exercice}
Simuler une file M/M/1 avec $\lambda = 3$ et $\mu = 5$ sur un horizon
de temps $T = 1000$. Estimer le nombre moyen de clients dans le
syst\`eme et comparer avec la valeur th\'eorique $L = 3/2$.
\end{exercice}

\begin{exercice}
Impl\'ementer l'algorithme de Gillespie pour le mod\`ele de
Lotka--Volterra stochastique~:
$X \xrightarrow{\alpha} 2X$, \quad $X + Y \xrightarrow{\beta} 2Y$,
\quad $Y \xrightarrow{\gamma} \emptyset$.
Comparer les trajectoires stochastiques au mod\`ele d\'eterministe.
\end{exercice}
